\section{The model}
\label{sec:model}

\subsection{Outcomes and the link function}

The unit of data is an outcome triple $(\rider, \metric, y)$, optionally indexed
by a site $\site$: rider $\rider$ experienced a condition summarised by the metric
value $\metric$ (for example, wind speed) at site $\site$ and produced a binary
outcome $y \in \{0,1\}$---a good session versus not, or a go versus no-go. We
model the outcome probability with a two-parameter latent factorisation,
\begin{equation}
  \label{eq:model}
  P(y{=}1 \mid \rider, \metric, \site)
    = \sigmoid\!\Big( a(\metric,\site)\,\big(\skill_\rider - \diff(\metric,\site)\big) \Big),
  \qquad \sigmoid(z) = \frac{1}{1+e^{-z}},
\end{equation}
with three latent objects:
\begin{itemize}
  \item $\skill_\rider \in \mathbb{R}$ --- the \emph{latent skill} of rider
    $\rider$, one scalar per rider (per pseudonym), identified up to a
    location/scale convention (Section~\ref{sec:ident});
  \item $\diff(\metric,\site)$ --- the \emph{latent difficulty} of the condition:
    how objectively hard the metric value $\metric$ is at site $\site$, a smooth
    function of $\metric$;
  \item $a(\metric,\site) > 0$ --- the \emph{discrimination}: how sharply the
    condition separates riders of different skill. In the Phase-1 form used
    throughout this paper, $a$ is a single scalar.
\end{itemize}
Equation~\eqref{eq:model} is the two-parameter logistic IRT model with the item
difficulty $b_i$ replaced by a function $\diff(\metric,\site)$ of a continuous
driver. Higher skill or lower difficulty raises the success probability; the two
enter only through their difference $\skill_\rider - \diff(\metric,\site)$, a fact
with consequences for identification (Section~\ref{sec:ident}).

\subsection{Difficulty as a continuous item}

Classical IRT difficulty is a per-item scalar because test items are discrete.
Environmental conditions are not: the metric $\metric$ is a continuum, so we model
$\diff(\metric,\site)$ as a smooth function---a spline (or, with full Bayesian
inference, a Gaussian process) over $\metric$---rather than a collection of
independent per-condition constants. Two properties make this well-posed. First,
smoothness lets neighbouring conditions share statistical strength, so a condition
observed only a few times still receives a sensible difficulty. Second, the
difficulty function is \emph{anchored to the physics-derived expert curve as its
prior}, which serves two distinct purposes: it fixes the difficulty scale in
interpretable, physical units, and it acts as a shrinkage prior toward which
sparsely observed conditions regress. This anchoring is about \emph{units and
regularisation}, not identification: $\diff$ is already identified once the skill
distribution is standardised (Section~\ref{sec:ident}); without the anchor it
would still be identified, but expressed only on an arbitrary scale.

\subsection{Recovering the single suitability curve as a marginal}
\label{sec:marginal}

The model in \eqref{eq:model} is strictly more general than today's single
suitability curve, and reduces to it exactly. Let $p(\skill)$ be the population
skill distribution. The single suitability curve is the marginal of
\eqref{eq:model} over that population,
\begin{equation}
  \label{eq:marginal}
  \mathrm{curve}(\metric)
    = \mathbb{E}_{\skill}\!\left[\, \sigmoid\!\big(a(\metric)(\skill - \diff(\metric))\big)\,\right]
    = \int \sigmoid\!\big(a(\metric)(\skill - \diff(\metric))\big)\, p(\skill)\, d\skill .
\end{equation}
Integrating skill out recovers precisely the population-level curve that a
conventional calibrator would fit. This is not merely an asymptotic
correspondence; in our implementation it is a regression test---the fitted
model's marginal is checked against a single-curve fit on the same data
(Section~\ref{sec:recovery}, the \emph{marginal matches single curve} gate).
Consequently, adopting Inverse Suitability adds a dimension to an existing
scoring engine; it does not fork it. Where a cell lacks the data to identify the
latent structure, the model falls back to \eqref{eq:marginal}, which is the
behaviour the engine already has.

\subsection{Skill-conditioned scoring}

For a chosen skill level $\skill$---for instance a population quantile mapped to
beginner, intermediate, or expert---the skill-conditioned suitability curve is
$\metric \mapsto \sigmoid\!\big(a(\metric)(\skill - \diff(\metric))\big)$. This is
the per-level scoring surface: the same condition is scored differently for a
beginner and an expert, while the underlying difficulty $\diff$ is invariant
(Figure~\ref{fig:atlas}). The skill level is treated as an explicit input, not
inferred from identity---a boundary we return to in Section~\ref{sec:governance}.
